% ===================================================================
%  সারণি নং 1 — বিদ্যালয়। — মাধ্যমিক বিদ্যালয়। — শ্রেণিকক্ষ।
%
%  Ceci n'est PAS une transcription : c'est une traduction, faite en
%  2026, en bengali d'aujourd'hui. D'ou trois differences avec
%  texto/io et texto/fr, et trois seulement :
%
%   * pas de \nl ni de \cc — la traduction ne suit aucune ligne
%     imprimee, et les lignes de ce fichier ne sont que du confort ;
%   * pas de \begin{VUpage} — elle n'a ni feuillet ni folio, et la
%     page de lecture ne lui en affiche pas ;
%   * le gras se pose sur le TERME seul, ponctuation dehors.
%
%  Tout le reste est identique, et doit l'etre : les cles « %%K » sont
%  celles de texto/io, macro pour macro pour l'apparat, et les renvois
%  \textsuperscript{(n)} visent les MEMES objets de la planche — ce
%  sont eux qui ouvrent les gros plans.
%
%  L'ido est le texte de base ; le francais sert de controle, et
%  l'anglais, l'espagnol et le hindi de calibrage.
%
%  LE GRAS NE SE POSE PAS SUR UN MOT-OUTIL. Le fac-simile ido le fait
%  par accident — « \VUgras{Ica} docisto » — et la page de lecture l'en
%  corrige ; la traduction, elle, n'a pas a repeter l'accident.
%
%  LE NOM DE L'EDITEUR SE PRONONCE « del-MASS ». Delmas est un nom du
%  Midi, et le s final s'y entend, comme dans Chaban-Delmas : on ecrit
%  donc দেলমাস, non দেলমা.
%
%  LES NOMS DE PERSONNE SONT NATURALISES, comme dans les huit autres
%  colonnes : Ioannes est জন, Iakobus est জেমস। Seuls le chien Medor
%  et l'ane Aliboron gardent leur nom.
%
%  LE BENGALI POSTPOSE SES RAPPORTS et place le verbe a la fin, comme
%  le hindi : on garde l'ordre des RENVOIS de l'ido, et l'on refait la
%  phrase autour.
%
%  LES CHIFFRES RESTENT ARABES, renvois et numeros d'alinea : ils
%  numerotent la gravure et les paragraphes, et doivent se lire d'une
%  colonne a l'autre. C'est le parti du hindi, et des sept autres.
% ===================================================================

%%K t01-apar-0 apar
\VUsaut{2.0mm}
\VUcentre{12.6pt}{120}{ব্যাখ্যা-পুস্তিকা}
\VUsaut{3.2mm}
\VUcentre{8.4pt}{160}{সংলগ্ন}
\VUsaut{2.6mm}
\VUtitre{15.6pt}{0}{দেলমাস সহায়ক সারণিসমূহের সঙ্গে}
\VUsaut{3.4mm}
\VUornamento{9pt}
\VUsaut{5.0mm}
\VUcentre{13.2pt}{90}{প্রথম পর্যায়}
\VUsaut{3.0mm}
\VUfilet{16mm}
\VUsaut{5.6mm}

%%K t01-apar-1 apar
\VUtitre{15.0pt}{60}{সারণি নং 1}
\VUsaut{1.4mm}
\VUtitre{11.4pt}{0}{বিদ্যালয়। --- মাধ্যমিক বিদ্যালয়। --- শ্রেণিকক্ষ।}
\VUsaut{2.2mm}
\VUfilet{20mm}
\VUsaut{6.4mm}

%%K t01-tit-1 sub
\VUpk{10.2pt}{40}{সাধারণ বর্ণনা।}
\VUsaut{3.0mm}

%%K t01-01-1 p
1. --- এই সারণিতে একটি মাধ্যমিক বিদ্যালয়ের \VUgras{শ্রেণিকক্ষ} দেখানো
হয়েছে। এই কক্ষে আটটি \VUgras{টেবিল} \textsuperscript{(18)} ও আটটি
\VUgras{বেঞ্চ} \textsuperscript{(32)} আছে। প্রতিটি বেঞ্চে তিনজন ছাত্র বসে।
\VUgras{শিক্ষক} \textsuperscript{(7)} একটি \VUgras{চেয়ারে} \textsuperscript{(8)}
বসে আছেন — চেয়ারটি তাঁর \VUgras{মঞ্চের} \textsuperscript{(6)} উপরে।

%%K t01-02-1 p
2. --- মঞ্চের বাঁ দিকে কাচের দরজাওয়ালা একটি \VUgras{আলমারি}
\textsuperscript{(15)} আছে। ডান দিকে \VUgras{ব্ল্যাকবোর্ড} \textsuperscript{(1)},
তার সঙ্গে \VUgras{ডাস্টার} \textsuperscript{(2)} ও \VUgras{চক}
\textsuperscript{(3)}।

%%K t01-02-2 p
মঞ্চের উপরে কয়েকটি \VUgras{দেয়ালচিত্র} \textsuperscript{(9, 11, 12)} ও একটি
\VUgras{মানচিত্র} \textsuperscript{(10)} ঝুলছে।

%%K t01-03-1 p
3. --- সব ছাত্র নিজের \VUgras{জায়গায়} \textsuperscript{(17)} নেই। জন
\textsuperscript{(4)} ব্ল্যাকবোর্ডের পাশে দাঁড়িয়ে আছে; ডাস্টার হাতে নিয়ে
\VUgras{জ্যামিতির চিত্রগুলি} মুছছে।

%%K t01-03-2 p
পিটার মঞ্চের কাছে; \VUgras{লিখিত অনুশীলনীগুলি} \textsuperscript{(5)} জড়ো করে
শিক্ষকের কাছে নিয়ে যাচ্ছে।

%%K t01-04-1 p
4. --- \VUgras{এই} শিক্ষক \VUgras{আধুনিক ভাষার} শিক্ষক। তিনি ছাত্রদের
বিদেশি ভাষা \VUgras{(জার্মান, ইংরেজি, স্প্যানিশ, ইতালীয়, রুশ} বা
\VUgras{ফরাসি)} বলতে, পড়তে ও লিখতে শেখান।

%%K t01-05-1 p
5. --- শ্রেণিকক্ষটি খুব প্রশস্ত ও খুব আলোকিত। শেষ দিকে দুটি
\VUgras{পাল্লাওয়ালা} \textsuperscript{(78)} একটি চওড়া \VUgras{জানালা} এবং হাওয়া
ও আলো আসার জন্য একটি \VUgras{কাচের দরজা} আছে।

%%K t01-05-2 p
এই কক্ষ ঠান্ডা নয়। ঘর গরম রাখার জন্য মাঝখানে একটি খুব ভালো
\VUgras{উনুন} \textsuperscript{(46)} আছে, তার সঙ্গে তার \VUgras{নল}
\textsuperscript{(45)}।

%%K t01-05-3 p
পার্টিশনে কয়েকটি \VUgras{আঁকশি} \textsuperscript{(58)} লাগানো আছে; সেগুলিতে
ছাত্রদের টুপি ও কোট ঝুলছে।

%%K t01-tit-2 sub
\VUsaut{5.2mm}
\VUpk{10.2pt}{40}{খেলার মাঠ।}
\VUsaut{3.6mm}

%%K t01-06-1 p
6. --- \VUgras{ঘড়িতে} \textsuperscript{(63)} নটা বেজে পাঁচ মিনিট।

%%K t01-06-2 p
\VUgras{বিরতি} \textsuperscript{(76)} এইমাত্র শেষ হল, আবার ক্লাস শুরু হচ্ছে।
বিরতি হয় \VUgras{উঠানে} \textsuperscript{(75)}। কয়েকজন ছাত্রকে খেলতে দেখা
যাচ্ছে। একজন \VUgras{তত্ত্বাবধায়ক} \textsuperscript{(74)} তাদের দেখাশোনা করছেন।

%%K t01-07-1 p
7. --- উঠানে আরও কয়েকজনকে দেখা যায়: \VUgras{পরিদর্শক}
\textsuperscript{(73)}, \VUgras{অধ্যক্ষ} \textsuperscript{(71)} ও \VUgras{পাঠ-তত্ত্বাবধায়ক}
\textsuperscript{(72)}।

%%K t01-07-2 p
পরিদর্শক বিদ্যালয়গুলি পরিদর্শন করেন, অধ্যক্ষ বিদ্যালয় পরিচালনা করেন,
পাঠ-তত্ত্বাবধায়ক পড়াশোনার তদারক করেন।

%%K t01-07-3 p
চতুর্থ ব্যক্তিটি \VUgras{দারোয়ান} \textsuperscript{(77)}। অনুপস্থিতদের নাম
লেখার জন্য বগলে একটি খাতা নিয়ে আছেন।

%%K t01-08-1 p
8. --- উঠানের শেষে \VUgras{ব্যায়ামের সরঞ্জাম} \textsuperscript{(70)} ও
\VUgras{হাতের কাজের} \textsuperscript{(69)} কর্মশালা আছে।

%%K t01-08-2 p
সারণির ডান দিকে \VUgras{সংগীতের} ও \VUgras{অঙ্কনের} \VUgras{কক্ষগুলি} চোখে
পড়ে।

%%K t01-noto-1 noto
\VUnotes{143.2mm}{%
(*) \textit{নামের} ক্ষেত্রে সেগুলিকে লাতিন রূপেও লিপিবদ্ধ করার পরামর্শ
দেওয়া হয়। অগণিত রূপভেদ এড়ানোর এটিই একমাত্র উপায়। তা ছাড়া,
\textit{Giovanni, Jean, Johann, Jan, Hans,} \textit{John, Ivan} প্রভৃতির
মধ্যে কোনটি বাছা হবে? নামের জন্য কি আলাদা অভিধান বানাতে হবে? লাতিন থেকে
তার কর্তৃকারক নিয়ে আমরা বলি: \VUgras{Ioannes, Ioanna; Iulius, Iulia;
Iakobus; Andreas; Lukas.} (সম্পূর্ণ ব্যাকরণ)।%
}

%%K t01-tit-3 sub
\VUpk{10.2pt}{40}{বিশদ বর্ণনা।}
\VUsaut{2.4mm}

%%K t01-tit-4 sub
\VUpk{10.2pt}{40}{ভালো ছাত্রেরা।}
\VUsaut{3.0mm}

%%K t01-09-1 p
9. --- মঞ্চের ডান দিকের প্রথম বেঞ্চের ছাত্রেরা হল \VUgras{হেনরি} (*),
\VUgras{স্টিফেন} ও \VUgras{চার্লস}।

%%K t01-09-2 p
হেনরি খুব মনোযোগী ও পরিশ্রমী; শিক্ষকের কথা শুনছে। তার টেবিলে একটি
\VUgras{খাতা} \textsuperscript{(28)}, একটি \VUgras{বই} \textsuperscript{(29)}, একটি
\VUgras{অভিধান} \textsuperscript{(30)} ও একটি \VUgras{কলমদানি}
\textsuperscript{(31)} আছে। তার বইয়ে অনেক \VUgras{পৃষ্ঠা}; প্রতিটি অধ্যায়ে
কয়েকটি \VUgras{অনুচ্ছেদ} ও প্রতিটি \VUgras{পঙ্‌ক্তিতে} অনেক
\VUgras{শব্দ}।

%%K t01-09-4 p
তার পাশের স্টিফেন \textsuperscript{(24)} অধ্যবসায়ী। পরের দিনের পড়া নিজের
খাতায় টুকে রাখছে। তার \VUgras{হাতের লেখা} ভালো। তার বই বন্ধ। কেবল
\VUgras{মলাটটি} \textsuperscript{(27)} দেখা যাচ্ছে। তার পাশেই তার
\VUgras{কম্পাস} \textsuperscript{(26)}।

%%K t01-09-5 p
চার্লস \textsuperscript{(23)} হাতে বই ধরে আছে; পড়া ঝালিয়ে নিতে বইটি খুলছে।
সব ছাত্রের মতো তার সামনেও একটি \VUgras{দোয়াত} \textsuperscript{(25)} আছে। সে
বাধ্য।

%%K t01-10-1 p
10. --- পাশের টেবিলে আছে \VUgras{লুই, অ্যান্টনি} ও \VUgras{পল}। লুই তার
\VUgras{পড়া} \textsuperscript{(22)} মুখস্থ বলছে এবং শিক্ষকের \VUgras{প্রশ্নের}
উত্তর দিচ্ছে। অ্যান্টনি \textsuperscript{(21)} খুব অন্যমনস্ক; শিক্ষক কখনও কখনও
তাকে শাস্তিও দেন। পল \textsuperscript{(20)} ভালো সঙ্গী, বিনয়ী ও উপকারী। সে খুব
মনোযোগী।

%%K t01-tit-5 sub
\VUsaut{4.2mm}
\VUpk{10.2pt}{40}{মন্দ ছাত্রেরা।}
\VUsaut{3.2mm}

%%K t01-11-1 p
11. --- হেনরির পিছনে \VUgras{জর্জ} \textsuperscript{(38)}। সে খারাপ ছেলে নয়,
কিন্তু \VUgras{বাচাল}, অন্যমনস্ক ও অস্থির। তার \VUgras{চপলতা} ও
\VUgras{অবাধ্যতা} ক্লাসে \VUgras{বিশৃঙ্খলা} ঘটায়, আর তাই প্রায়ই সে কড়া
\VUgras{তিরস্কার} পাওয়ার যোগ্য হয়। তার \VUgras{খসড়া খাতা}
\textsuperscript{(35)} নোংরা, \VUgras{কলম} \textsuperscript{(34)} দিয়ে সে সেটিকে
\VUgras{কালির দাগে} ভরে ফেলে, আর তাই সব সময় তার \VUgras{রবার}
\textsuperscript{(36)} ব্যবহার করে; তার \VUgras{ব্যাগ} \textsuperscript{(33)} ছেঁড়া,
কারণ সে বড় অসাবধান। আধুনিক ভাষার ক্লাসে সে কেন তার \VUgras{পেনসিলদানি}
\textsuperscript{(37)} আনে?

%%K t01-11-2 p
তার পাশের এডমন্ড \textsuperscript{(39)} তাকে পছন্দ করে না। সত্যি বলতে সে কিছুটা
অলস, কিন্তু চুপচাপ ও শান্ত। সে বকবক করে না; কেবল মন দেওয়ার বদলে
\VUgras{পাঠ্যবইয়ের} \VUgras{গল্পগুলি} পড়তে বেশি ভালোবাসে।

%%K t01-11-3 p
এই বেঞ্চের তৃতীয় ছাত্র \VUgras{লুদোভিক} \textsuperscript{(40)}। সে
অধ্যবসায়ী। সাদা এক \VUgras{কাগজের পাতায়} লিখছে। তার সামনে রাখা আছে তার
\VUgras{শুষ্ক কাগজ} \textsuperscript{(41)}।

%%K t01-12-1 p
12. --- শিক্ষকের বাঁ দিকের দ্বিতীয় বেঞ্চের ছাত্রেরা খুব ছোট। প্রথমজন,
ছোট্ট \VUgras{এমিল}, এখনও ভালো করে লিখতে জানে না; সে তার \VUgras{স্লেট}
\textsuperscript{(42)}, তার \VUgras{স্লেট-পেনসিল} ও তার \VUgras{স্পঞ্জ}
\textsuperscript{(43)} নিয়ে আসে।

%%K t01-12-2 p
তার পাশে \VUgras{অগাস্ট} ও \VUgras{বার্নার্ড} \textsuperscript{(44)}। শেষজন ভালো
কাজ করে, কিন্তু তার \VUgras{বেশভূষা} বড় অগোছালো।

%%K t01-13-1 p
13. --- তাদের পিছনে তিনজন \VUgras{অলস}। \VUgras{আলবার্ট} পড়া মুখস্থ করে
না। \VUgras{জেমস} \textsuperscript{(47)} শোনার বদলে ঘুমায়। তার \VUgras{আলস্য}
তার জন্য \VUgras{ভর্ৎসনা} এমনকি \VUgras{শাস্তিও} ডেকে আনে। শেষে,
\VUgras{ক্রিশ্চিয়ান} \textsuperscript{(48)} বড় বেশি লঘুচিত্ত। সে খারাপ
\VUgras{আচরণ} করে এবং শোধরানোর কোনো চেষ্টাই করে না।

%%K t01-14-1 p
14. --- তার পাশের ছাত্রেরা অবশ্য ভালো আচরণ করে। \VUgras{আর্নেস্ট}
\textsuperscript{(51)} শৃঙ্খলা ও নিয়ম ভালোবাসে; সব সময় তার \VUgras{স্কেল}
\textsuperscript{(50)} ও তার \VUgras{পেনসিল} \textsuperscript{(49)} ব্যবহার করে। সে
\VUgras{দিবা-ছাত্র}।

%%K t01-14-2 p
\VUgras{ফ্রান্সিস} \textsuperscript{(52)} \VUgras{আবাসিক ছাত্র}; সে বিদ্যালয়ের
\VUgras{পোশাক} পরে; সেখানেই খায়, সেখানেই ঘুমায়। সে ভালো \VUgras{সঙ্গী}।

%%K t01-14-3 p
মরিস তার \VUgras{ছুরি} \textsuperscript{(56)} দিয়ে তার \VUgras{পেনসিল} ছাঁটছে;
সে খুব যত্নবান।

%%K t01-14-4 p
সে তার সব বই আনে, তার \VUgras{ব্যাকরণ} \textsuperscript{(55)}, তার
\VUgras{নোটবই} \textsuperscript{(54)} এমনকি একটি \VUgras{লেখার তক্তাও}
\textsuperscript{(53)}। তার \VUgras{স্কুলব্যাগ} \textsuperscript{(57)} তার পিছনে
মেঝেতে পড়ে আছে।

%%K t01-14-5 p
কক্ষের শেষ দিকের দুটি টেবিল \VUgras{ভালো ছাত্রেরা} \textsuperscript{(19)}
দখল করে আছে।

%%K t01-tit-6 sub
\VUsaut{4.6mm}
\VUpk{10.2pt}{40}{অঙ্কন ও সংগীত।}
\VUsaut{3.4mm}

%%K t01-15-1 p
15. --- \VUgras{অঙ্কনকক্ষে} দেয়ালে সুন্দর \VUgras{নমুনা} ঝুলছে এবং ছাদ
থেকে \VUgras{ঢাকনিসহ} একটি \VUgras{বাতি} \textsuperscript{(62)} ঝুলছে।
\VUgras{শিক্ষক} \textsuperscript{(60)} খুব ধৈর্যশীল; ছাত্রদের \VUgras{অঙ্কন}
\textsuperscript{(61)} শুধরে দেন এবং সরাসরি দেখে একটি \VUgras{আবক্ষ মূর্তি}
\textsuperscript{(59)} আঁকতে শেখান।

%%K t01-16-1 p
16. --- \VUgras{সংগীতকক্ষটি} খুব কৌতূহলজনক মনে হয়। সেখানে
\VUgras{সংগীত} পড়তে শেখা যায়, \VUgras{স্বরলিপির পাতায়} লেখা
\VUgras{সপ্তকের স্বরগুলি} \textsuperscript{(67)} (সা রে গা মা পা ধা নি\,=\,C.
D. E. F. G. A. B.) গাইতে শেখা যায়, \VUgras{ক্লেফ} চিনতে এবং সুন্দর
\VUgras{সুর} গাইতে শেখা যায়। স্বরলিপি রাখা হয় \VUgras{স্ট্যান্ডে}
\textsuperscript{(65)}। একজন ছাত্র \VUgras{পিয়ানো বাজাচ্ছে} \textsuperscript{(64)}
এবং \VUgras{সংগীতশিক্ষক} \textsuperscript{(66)} \VUgras{তাল দিচ্ছেন}
\textsuperscript{(68)}।

%%K t01-17-1 p
17. --- \VUgras{হাতের কাজের} \VUgras{কর্মশালার} \textsuperscript{(69)} এক কোণ
চোখে পড়ে, যেখানে ছেলেরা কাঠ রাঁদা করতে ও লোহা রেতে শিখতে পারে। সব
মাধ্যমিক বিদ্যালয়ে এটি নেই।

%%K t01-tit-7 sub
\VUsaut{5.0mm}
\VUpk{10.2pt}{40}{বিজ্ঞানকক্ষ।}
\VUsaut{3.6mm}

%%K t01-18-1 p
18. --- গণিতের শিক্ষক ছাত্রদের \VUgras{জ্যামিতি,} \VUgras{পাটিগণিত,
হিসাব} ও \VUgras{মেট্রিক পদ্ধতি} শেখান। ব্ল্যাকবোর্ডে জ্যামিতির প্রধান
চিত্রগুলি দেখা যাচ্ছে: একটি \VUgras{বৃত্ত} \textsuperscript{(a)}, একটি
\VUgras{বর্গ} \textsuperscript{(b)}, একটি \VUgras{ত্রিভুজ} \textsuperscript{(c)},
একটি \VUgras{রেখা} \textsuperscript{(d)}।

%%K t01-18-2 p
একটি \VUgras{অঙ্কও} দেখা যাচ্ছে; এটি \VUgras{যোগ} নয়, \VUgras{বিয়োগও}
নয়, \VUgras{গুণও} নয়: এটি \VUgras{ভাগ} \textsuperscript{(e)}।

%%K t01-19-1 p
19. --- মেট্রিক পদ্ধতির চিত্রে চোখে পড়ে: একটি \VUgras{মিটার}
\textsuperscript{(a)}, একটি \VUgras{দাঁড়িপাল্লা} \textsuperscript{(b)} ও তার
\VUgras{বাটখারা}, একটি \VUgras{লিটার} \textsuperscript{(e)}, একটি
\VUgras{ডেকালিটার} \textsuperscript{(f)}, একটি \VUgras{ডেসিলিটার} ও একটি
\VUgras{ঘনমিটার} \textsuperscript{(d)}।

%%K t01-19-2 p
ছাত্রেরা \VUgras{প্রাকৃতিক বিজ্ঞানও} \textsuperscript{(11)} পড়ে ---
প্রাণিবিদ্যা \textsuperscript{(a)}, উদ্ভিদবিদ্যা \textsuperscript{(b)},
ভূতত্ত্ব \textsuperscript{(c)} --- এবং \VUgras{ইতিহাস} \textsuperscript{(12)}।

%%K t01-19-3 p
আলমারিতে \VUgras{পদার্থবিদ্যার} \textsuperscript{(15)} ও \VUgras{রসায়নের}
\textsuperscript{(16)} যন্ত্রপাতি আছে।

%%K t01-19-4 p
আলমারির উপরে আছে: একটি \VUgras{গোলক} \textsuperscript{(14)}, একটি
\VUgras{শঙ্কু} ও একটি \VUgras{ভূগোলক} \textsuperscript{(13)}।

%%K t01-19-5 p
চিত্রগুলির উপরে একটি \VUgras{মানচিত্র} ঝুলছে, যেখানে পৃথিবীর পাঁচটি অংশ
দেখানো হয়েছে: \VUgras{ইউরোপ} \textsuperscript{(g)}, \VUgras{এশিয়া}
\textsuperscript{(e)}, \VUgras{আফ্রিকা} \textsuperscript{(f)}, \VUgras{আমেরিকা}
\textsuperscript{(ab)} ও \VUgras{ওশেনিয়া} \textsuperscript{(h)}, আর দুটি বড়
মহাসাগর, \VUgras{প্রশান্ত} \textsuperscript{(c)} ও \VUgras{আটলান্টিক}
\textsuperscript{(d)}।
\VUsaut{6.0mm}
\VUfilet{22mm}
